% --- Suggested replacement/addition for Section VII-B ---
A constraint violation was counted when an accelerating action was executed while the true simulated inter-vehicle gap was below
\(d_{\mathrm{safe}}=v t_{\mathrm{reaction}}+d_{\min}\).
The runtime safety gate, however, operated only on the noisy edge-estimated gap.
This separates the outcome criterion from the gate predicate and permits gate errors under sensor uncertainty to remain observable.

To isolate component effects, we additionally evaluated four ablations under the same 100 seeds and disturbance schedules:
freshness checking only, local fallback only, runtime safety gating only, and the full configuration combining all three.
We report true-gap violation rate, stale-cloud-action rate, blocked-action rate, MTTR, minimum true gap, mean ego speed,
mean absolute jerk, and non-recovered episodes. Trace completeness is retained as an auditability property but is not interpreted
as an independent performance effect because the required artifact set is configured by design.

% --- Suggested interpretation paragraph ---
The ablation separates complementary mechanism effects. Freshness checking eliminates stale-cloud execution in this bounded model,
local fallback determines immediate recovery from cloud outages, and the safety gate reduces unsafe acceleration relative to a
ground-truth outcome criterion that is distinct from the noisy state estimate used by the gate itself. The full configuration
combines these effects but also incurs behavioral costs visible in minimum-gap, speed, and jerk measures. Accordingly, the mobility
illustration supports executability and metric sensitivity within the declared model rather than general safety effectiveness.

% --- Suggested Section VII-C sentence ---
True-state constraint violations and blocked candidate actions provide complementary evidence for P7: the former measures safety
outcomes against the simulated ground truth, whereas the latter measures gate intervention frequency.
